\subsection{\texorpdfstring{\(\mathfrak{sl}(2,k)\) 的典则基}{\textbackslash mathfrak\{sl\}(2,k) 的典则基}}

引理 1. 设 \(\mathbf{A}\) 是 \(k\) 上的结合代数，\(H\) 与 \(X\) 是 \(\mathbf{A}\) 中满足 \([H, X] = 2X\) 的元素。

\begin{enumerate}
\def\labelenumi{(\roman{enumi})}
\item
  \([H,X^n ] = 2nX^n\) 对任何整数 \(n\geq 0\) 成立
\item
  若 \(Z\) 是 \(A\) 中满足 \([Z, X] = H\) 的元素，则对任何整数 \(n > 0\) ，
\end{enumerate}

\[
[ Z, X ^ {n} ] = n X ^ {n - 1} (H + n - 1) = n (H - n + 1) X ^ {n - 1}.
\]

映射 \(T \mapsto [H, T]\) （从 \(A\) 到 \(A\) 的）是一个导子，由此得 (i)。在 (ii) 的假定下，

\[
\begin{array}{l} [ Z, X ^ {n} ] = \sum_ {i + j = n - 1} X ^ {i} H X ^ {j} \\ \qquad = \sum_ {i + j = n - 1} (X ^ {i} X ^ {j} H + X ^ {i} 2 j X ^ {j}) \\ \qquad = n X ^ {n - 1} H + 2 X ^ {n - 1} \frac {n (n - 1)}{2} \\ \qquad = n X ^ {n - 1} (H + n - 1). \end{array}
\]

另一方面，由 (i) 有 \(X^{n-1}(H+n-1)=(H-n+1)X^{n-1}\) 。证毕。

回顾我们用 \(\mathfrak{sl}(2,k)\) 表示由 2 阶、迹为零、元素在 k 中的方阵组成的李代数。此李代数是 3 维单李代数（第一章 §6 第 7 节 例）。\(\mathfrak{sl}(2,k)\) 的典则基是基 \((X_{+},X_{-},H)\) ，其中

\[
X _ {+} = \left( \begin{array}{c c} 0 & 1 \\ 0 & 0 \end{array} \right) \quad X _ {-} = \left( \begin{array}{c c} 0 & 0 \\ - 1 & 0 \end{array} \right) \quad H = \left( \begin{array}{c c} 1 & 0 \\ 0 & - 1 \end{array} \right).
\]

我们有

\[
[ H, X _ {+} ] = 2 X _ {+} \quad [ H, X _ {-} ] = - 2 X _ {-} \quad [ X _ {+}, X _ {-} ] = - H.\tag{1}
\]

由于 \(\mathfrak{sl}(2,k)\) 的恒同表示是单射，H 是 \(\mathfrak{sl}(2,k)\) 的半单元素，而 \(X_{+}, X_{-}\) 是 \(\mathfrak{sl}(2,k)\) 的幂零元素（第一章 §6 第 3 节 定理 3）。由第七章 §2 第 1 节 例 4，kH 是 \(\mathfrak{sl}(2,k)\) 的嘉当子代数。映射 \(U \mapsto -^{t}U\) 是李代数 \(\mathfrak{sl}(2,k)\) 的对合自同构，称为 \(\mathfrak{sl}(2,k)\) 的典则对合；它把 \((X_{+}, X_{-}, H)\) 变为 \((X_{-}, X_{+}, -H)\) 。

引理 2. 在 \(\mathfrak{sl}(2,k)\) 的包络代数中，

\[
[ H, X _ {+} ^ {n} ] = 2 n X _ {+} ^ {n} \quad [ H, X _ {-} ^ {n} ] = - 2 n X _ {-} ^ {n}
\]

对任何整数 \(n \geq 0\) 成立，且

\[
[ X _ {-}, X _ {+} ^ {n} ] = n X _ {+} ^ {n - 1} (H + n - 1) = n (H - n + 1) X _ {+} ^ {n - 1}
\]

\[
[ X _ {+}, X _ {-} ^ {n} ] = n X _ {-} ^ {n - 1} (- H + n - 1) = n (- H - n + 1) X _ {-} ^ {n - 1}
\]

若 n \textgreater{} 0。

第一与第三个关系式由引理 1 得出。其余的可用 \(\mathfrak{sl}(2,k)\) 的典则对合从它们推出。

